\section{LZSS Encoder/Decoder}

\subsection{Purpose}

LZSS is the core compression algorithm.
It scans a 1-D byte stream and replaces repeated sequences with compact
back-references into a sliding history window, emitting only the tokens and flag
bits needed by the decoder to reconstruct the original stream.

All other modules (scanner, model) are pre-processors that rearrange or transform
the data to increase the regularity that LZSS then exploits.

\subsection{Files}

\begin{itemize}
  \item \texttt{include/lzss.hpp} --- constants and public API declarations.
  \item \texttt{src/lzss.cpp} --- private data structures and implementations
        (\textasciitilde{}230 lines).
\end{itemize}

\subsection{Constants}

\begin{center}
\begin{tabular}{>{\ttfamily}l r l}
\toprule
\textbf{Constant} & \textbf{Value} & \textbf{Rationale} \\
\midrule
LZSS\_WINDOW\_SIZE & 4096 & 12-bit offset; one row of a 256-wide image = 256 B, 16 rows fit \\
LZSS\_LOOKAHEAD    & 32   & 5-bit length field (values 0..31 = actual lengths 3..34 bytes) \\
LZSS\_MIN\_MATCH   & 3    & Match token costs 18 bits; 3 literals cost 27 bits --- net saving at 3 \\
LZSS\_OFFSET\_BITS & 12   & Token: 12 bits for the window offset \\
LZSS\_LENGTH\_BITS & 5    & Token: 5 bits for adjusted length (length $-$ MIN\_MATCH) \\
\bottomrule
\end{tabular}
\end{center}

\textbf{MIN\_MATCH derivation:}
A match token costs $1 + 12 + 5 = 18$ bits.
Three literals each cost $1 + 8 = 9$ bits, totalling 27 bits.
A match of length 3 therefore saves 9 bits.
A match of length 2 would cost the same as 2 literals (18 bits vs $2 \times 9 = 18$ bits)
--- no saving --- so 3 is the minimum meaningful match length.

\subsection{Internal Data Structures}

\subsubsection{Window}

\begin{lstlisting}[style=cpp]
struct Window {
    uint8_t  buf[LZSS_WINDOW_SIZE] = {};
    uint32_t pos = 0;

    void push(uint8_t byte) {
        buf[pos] = byte;
        pos = (pos + 1) % LZSS_WINDOW_SIZE;
    }
};
\end{lstlisting}

A circular byte buffer of 4096 bytes representing the sliding history.
\texttt{pos} always points to the slot that will receive the \emph{next} byte.
Both encoder and decoder maintain an identical \texttt{Window} instance so that
the decoder can reconstruct matches from its own copy of history without any
explicit dictionary transmission.

\subsubsection{FlagGroup}

\begin{lstlisting}[style=cpp]
struct FlagGroup {
    static constexpr int GROUP = 8;
    uint8_t  flags[8];
    uint32_t offsets[8];   // match: window offset
    uint32_t lengths[8];   // match: length - LZSS_MIN_MATCH (0-based)
    uint8_t  literals[8];  // literal byte
    int count = 0;         // tokens accumulated (0..8)

    void add_literal(uint8_t byte);
    void add_match(uint32_t offset, uint32_t length);
    bool full() const;
    void flush(BitWriter& out);
};
\end{lstlisting}

Collects up to 8 tokens before emitting them as a group.
When \texttt{flush()} is called:
\begin{enumerate}
  \item One \emph{flag byte} is emitted, with each bit (MSB first) indicating whether
        the corresponding token is a match (1) or a literal (0).
        If the group is partial (fewer than 8 tokens), the remaining bits are padded with 0.
  \item Each token's payload is emitted in order:
        literal $\to$ 8 bits; match $\to$ 12-bit offset + 5-bit adjusted length.
\end{enumerate}

The length stored in the token is \texttt{length} $-$ \texttt{LZSS\_MIN\_MATCH},
so a 3-byte match stores 0 and a 34-byte match stores 31.
The decoder adds \texttt{LZSS\_MIN\_MATCH} back during decoding.

\subsubsection{HashChains}

\begin{lstlisting}[style=cpp]
struct HashChains {
    static constexpr uint32_t HASH_SIZE = 65536;
    static constexpr uint32_t NIL = LZSS_WINDOW_SIZE;  // sentinel

    uint32_t head[65536];          // head[h] = most-recent window slot with hash h
    uint32_t prev[LZSS_WINDOW_SIZE]; // prev[slot] = previous slot with same hash

    static uint32_t hash2(uint8_t a, uint8_t b); // (a << 8) | b
    void insert(uint32_t slot, uint8_t a, uint8_t b);
    std::pair<uint32_t,uint32_t> find_match(...);
};
\end{lstlisting}

A hash-based match finder using linked lists within the sliding window.

\textbf{Hash function:} \texttt{hash2(a, b) = (a << 8) | b} produces a 16-bit key
from the first two bytes of a candidate position.
The 65536-bucket hash table maps each 2-byte pair to the most recent window slot
that starts with those two bytes.

\textbf{insert(slot, a, b):} prepends \texttt{slot} to the chain for hash \texttt{(a,b)}.
The existing head becomes \texttt{prev[slot]}.

\textbf{find\_match(pattern, len, win):} walks the chain for
\texttt{hash2(pattern[0], pattern[1])}, comparing up to \texttt{LZSS\_LOOKAHEAD}
bytes at each candidate position.
The chain walk is capped at 128 iterations (\texttt{chain\_limit}) to bound worst-case
encoder cost on pathological inputs.

Matches shorter than \texttt{LZSS\_MIN\_MATCH} are discarded; the function returns
\texttt{(0, 0)} in that case.
Otherwise it returns the \texttt{(offset, length)} pair for the longest match found.

\subsection{lzss\_encode()}

\begin{lstlisting}[style=cpp]
void lzss_encode(const std::vector<uint8_t>& input, BitWriter& out);
\end{lstlisting}

Delegates to the private \texttt{encode\_impl()} function, which:

\begin{enumerate}
  \item Initialises \texttt{Window}, \texttt{HashChains}, \texttt{FlagGroup}, position \texttt{pos=0}.
  \item While \texttt{pos < input.size()}:
    \begin{enumerate}
      \item Calls \texttt{HashChains::find\_match(input + pos, lookahead, win)}.
      \item If length $\geq$ \texttt{MIN\_MATCH}: \texttt{FlagGroup::add\_match}; advance \texttt{pos}
            by \texttt{length}; push each consumed byte to the window and insert its hash.
      \item Else: \texttt{FlagGroup::add\_literal(input[pos])}; advance by 1; push and insert.
      \item If \texttt{FlagGroup::full()}: \texttt{FlagGroup::flush(out)}.
    \end{enumerate}
  \item Flush the final (partial) group.
\end{enumerate}

The caller is responsible for calling \texttt{BitWriter::flush()} after
\texttt{lzss\_encode()} returns.

\subsection{lzss\_encode\_to\_buffer()}

\begin{lstlisting}[style=cpp]
std::vector<uint8_t> lzss_encode_to_buffer(const std::vector<uint8_t>& input);
\end{lstlisting}

Wraps \texttt{encode\_impl} with a \texttt{std::ostringstream} to collect output in memory.
After encoding, extracts the string into a \texttt{vector\textless uint8\_t\textgreater}.
Used by \texttt{codec.cpp} in adaptive mode to check compressed size before committing
to the output file (see Section~\ref{sec:codec-compress}).

\subsection{lzss\_decode()}

\begin{lstlisting}[style=cpp]
void lzss_decode(BitReader& in, std::vector<uint8_t>& output, size_t expected_size);
\end{lstlisting}

Reconstructs exactly \texttt{expected\_size} bytes from the compressed bit stream:

\begin{enumerate}
  \item Initialises \texttt{Window}, clears and reserves \texttt{output}.
  \item While \texttt{output.size() < expected\_size}:
    \begin{enumerate}
      \item Reads one 8-bit flag byte.
      \item For each bit $i$ from 7 down to 0 (MSB first), while bytes still needed:
        \begin{itemize}
          \item Bit = 0 (literal): read 8-bit value; append to \texttt{output}; push to window.
          \item Bit = 1 (match): read 12-bit offset, 5-bit adjusted length.
                Actual length = adjusted + \texttt{MIN\_MATCH}.
                Copy \texttt{length} bytes from the window at distance \texttt{offset},
                appending each byte to \texttt{output} and pushing it back to the window.
                The copy handles overlapping matches (where \texttt{length > offset})
                by clamping \texttt{effective\_k = k \% offset} --- this reproduces the
                run-length extension behaviour of the encoder.
        \end{itemize}
    \end{enumerate}
\end{enumerate}

The decoder stops as soon as \texttt{output.size() == expected\_size}, even if the
flag byte's remaining bits have not been consumed.
A partial flag group at the end of the stream is therefore normal and expected.
